\documentclass[a4paper, 12pt]{article}

\usepackage{utils}

\renewcommand*{\today}{26 March 2026}

\begin{document}

\hotbox{Analyse 4}{CM 13}{\today}

\begin{lemme}{Riemann-Lebesgue}{}
    Soit $f : ]a, b[ \to \mathbb{C}$ une fonction intégrable sur $]a, b[$. Alors

    $$
    \lim_{t \to +\infty} \int_a^b f(x) e^{i t x} \, dx = 0
    $$
\end{lemme}

\begin{demonstration}
    \begin{description}
        \item[1 $f$ est constante.]
        On a $f(x) = f_0$ pour tout $x \in ]a, b[$. Alors
        $$
        \int_a^b f(x) e^{i t x} \, dx = f_0 \int_a^b e^{i t x} \, dx = f_0 \frac{e^{i t b} - e^{i t a}}{i t}
        $$
        Ainsi on a
        $$
        \left| \int_a^b f(x) e^{i t x} \, dx \right| = \left| f_0 \frac{e^{i t b} - e^{i t a}}{i t} \right| \leq  \frac{|f_0|}{t} (|e^{i t a}| + |e^{i t b}|) \leq \frac{2 |f_0|}{t}
        $$
        et donc
        $$
        \lim_{t \to +\infty} \left| \int_a^b f(x) e^{i t x} \, dx \right| = 0
        $$

        \item[2 $f$ est une fonction en escalier.]
        Soit $f$ une fonction en escalier sur $]a, b[$. On peut trouver une subdivision $a = x_0 < x_1 < \cdots < x_n = b$ de $[a, b]$ telle que $f$ soit constante sur chaque intervalle $]x_{k-1}, x_k[$ pour $k = 1, \ldots, n$.

        On introduit la fonction d'Heaveside $H_{x_i} : \R \to \R$ définie par
        $$
        H_{x_i}(x) = \begin{cases}
        0 & \text{si } x < x_i \\
        1 & \text{si } x \geq x_i
        \end{cases}
        $$
        Ainsi on a
        $$
        H_{x_i} - H_{x_{i+1}} = \begin{cases}
            1 & \text{si } x \in [x_i, x_{i+1}[ \\
            0 & \text{sinon } \\
        \end{cases} 
        $$

        Par conséquent, on peut écrire $f$ comme
        $$
        f(x) = \sum_{j=0}^{n-1} f_j(H_{x_j}(x) - H_{x_{j+1}}(x))
        $$

        On a
        \begin{align*}
            \int_a^b f(x) e^{i t x} \, dx &= \int_a^b \sum_{j=0}^{n-1} f_j(H_{x_j}(x) - H_{x_{j+1}}(x)) e^{i t x} \, dx \\
            &= \sum_{j=0}^{n-1} \int_{x_j}^{x_{j+1}} f_j e^{i t x} \, dx
        \end{align*}
        Ainsi on a une somme de fonctions du type $f_j e^{i t x}$, et on peut appliquer le cas 1 à chacune de ces fonctions pour conclure que
        $$
        \lim_{t \to +\infty} \int_a^b f(x) e^{i t x} \, dx = 0
        $$
        \stump{Pas fini}{rohhh}

        \item[{3 $f$ est une fonction intégrable sur $]a, b[$.}]
        
        Par définition de l'intégrabilité, pour tout $\varepsilon > 0$, il existe une subdivision $a = x_0 < x_1 < \cdots < x_n = b$ de $[a, b]$
        et une fonction $\psi$ compatible avec la subdivision telle que
        $$
        \int_a^b |f(x) - \psi(x)| \, dx < \varepsilon
        $$
    \end{description}
\end{demonstration}

\end{document}